\documentclass[twocolumn]{aastex631}
\shorttitle{Little red dots should turn red over time}
\shortauthors{Pandev}

\begin{document}

\title{Little Red Dots Should Turn Red Over Cosmic Time: A Predicted
Subtype March and the Blue Dominance of the $z>8.5$ Population}

\author[0009-0005-8153-071X]{Stilian Pandev}
\affiliation{Independent researcher}
\email{stilianpandev@gmail.com}

\correspondingauthor{Stilian Pandev}

\begin{abstract}
Stacked spectroscopy of little red dots (LRDs) resolves an ordered
continuum-subtype sequence, from extreme red (xLRD) to blue (bLRD)
classes, that inverts to a monotonic sequence in the fraction of ionizing
photons processed by dense nuclear gas. If this spectral diversity is a
visibility sequence---one population seen at different stages of nuclear
envelope growth---rather than distinct formation routes, the subtype
mixture cannot be constant: it must evolve with redshift. We extract this
prediction, with uncertainties, from a semi-analytic lifecycle model in
which the envelope covering odds equal the squared engine-to-host
continuum contrast. Averaged over realizations, the selected red fraction
(xLRD+plusLRD) rises from $0.04\pm0.05$ at $z\simeq9.5$ through
$0.09\pm0.04$ and $0.23\pm0.05$ to $0.49\pm0.09$ at $z\simeq4.5$
(plateau-consistent $0.48\pm0.14$ at $z\simeq3.25$, an extrapolation
below the calibration range), while blue-continuum LRDs are the largest
class at $z>8.5$ ($0.61\pm0.15$ of the selected population), with xLRDs
essentially absent there. The march is population physics, not selection:
removing the survey selection strengthens it, and the rise to
$z\simeq4.5$ is monotonic in all realizations. It is testable now by
splitting existing stacked and spectroscopic samples (predicted slope
$0.10$--$0.14$ per unit redshift between $z\simeq7.5$ and 4.5). The same
physics implies that the reported $z\simeq9.5$ number density cannot
currently constrain formation: the intrinsic UV-bright nuclear-active
population exceeds it $\sim$34-fold but is blue and removed by the
redness selection. A sample of ten classified $z>8.5$ LRDs with a red
majority would exclude the prediction at $\ge99\%$ confidence even at its
upper uncertainty edge.
\end{abstract}

\keywords{Active galactic nuclei (16) --- Supermassive black holes (1663)
--- High-redshift galaxies (734)}

\section{Introduction}\label{sec:intro}
Little red dots---compact, red, frequently broad-line JWST sources at
$z\gtrsim4$ \citep{Matthee2024,Greene2024,Kocevski2024,Kocevski2025}---show
an ordered spectral-subtype sequence in stacked spectroscopy
\citep{PerezGonzalez2026}. Inverting a dense-plus-diffuse two-zone
photoionization model through the four stacks yields dense photon
fractions of 0.981, 0.901, 0.738, and 0.086 from the xLRD through bLRD
classes at nearly constant diffuse ionization parameter, with the
population-level He~II constraint \citep{Sok2026} selecting a soft
extreme-UV ionizing basis in that inversion. This ordered, one-parameter
structure is what a \emph{visibility sequence} looks like: a single
population whose nuclear envelope intercepts a varying share of the
ionizing output, rather than four unrelated formation routes.
Dense-envelope models of individual LRDs
\citep[e.g.,][]{Naidu2025,deGraaff2025} have argued on independent
grounds that envelope properties track the accretion state; spectral
diversity has been mapped to varying black-hole-to-host luminosity ratios
in a stationary framework \citep{Barro2025}; two LRDs have been observed
transitioning into quasars at the phase exit \citep{Fu2025}; and
low-spin halo models address the redshift evolution of LRD observability
\citep{PacucciLoeb2025}. The population-level consequence for the subtype
mixture, however, appears not to have been stated as a quantitative,
falsifiable prediction with a slope and uncertainties, which is the
purpose of this Letter: if the dense photon share is set by
the growth state of the nuclear envelope, and envelopes grow as nuclear
reservoirs feed while stellar hosts are still assembling, then the
subtype \emph{mixture must evolve with redshift}. A population seen
earlier in its lifecycle must be bluer.

\section{The model in brief}\label{sec:model}
The prediction is extracted from a semi-analytic lifecycle model
developed in a companion paper (Pandev 2026b, submitted). Halo histories,
importance-reweighted to the \citet{Tinker2008} mass function, carry
nuclear gas and black-hole reservoirs through an explicit five-state
lifecycle with competing exit hazards; the observation layer assigns each
source a porous nuclear reprocessor, a two-zone soft-EUV emission model
built on thermal-balance Cloudy anchors \citep{Chatzikos2023}, and the
published $L_{5100}/L_{2500}$ subtype boundaries (1.8, 3.1, 6.3). The
envelope covering $C$ follows an equilibrium law: the covering odds equal
the squared intrinsic engine-to-host continuum contrast,
\begin{equation}
\frac{C}{1-C}=x^{2},\qquad
x\equiv\frac{L_{\nu,\rm engine}(2500\,\text{\AA})}
{L_{\nu,\star}(2500\,\text{\AA})},
\label{eq:law}
\end{equation}
applied with a frozen per-halo lognormal scatter of 0.35 dex. In the
companion paper this law is selected against the stacked line ratios and
the $z\le7.5$ number densities, its exponent is constructed through an
explicit elimination program (momentum-limited envelope building against
drag-limited clump removal; three alternative mechanisms rejected by the
data), and the resulting photon partition matches the eight stacked
[O~III] ratios at least as well as an imposed closure while improving the
reliable-bin demographic fit to $\chi^{2}=1.2$. Here only one property
matters: covering---and therefore subtype---tracks the engine's contrast
with the growing host.

\section{The subtype march}\label{sec:march}

\begin{figure}
\centering
\includegraphics[width=\columnwidth]{subtype_march_with_bands.pdf}
\caption{The predicted subtype march: seed-averaged selected subtype
fractions versus redshift, with combined binomial and
realization-to-realization uncertainties. Subtypes follow the published
$L_{5100}/L_{2500}$ boundaries; weighting is by selection probability
times comoving weight. The $z\simeq3.25$ bin lies below the demographic
calibration range and is an extrapolation.}
\label{fig:march}
\end{figure}

\begin{deluxetable}{lccccc}
\tablecaption{Predicted selected subtype fractions (seed mean $\pm$
combined binomial and seed uncertainty). $N_{\rm eff}$ is the effective
selection-weighted sample per bin.\label{tab:march}}
\tablehead{\colhead{$z$} & \colhead{xLRD} & \colhead{plusLRD} &
\colhead{minusLRD} & \colhead{bLRD} & \colhead{$N_{\rm eff}$}}
\startdata
9.5 & $0.00\pm0.00$ & $0.04\pm0.05$ & $0.35\pm0.16$ & $0.61\pm0.15$ & 21\\
7.5 & $0.00\pm0.00$ & $0.09\pm0.04$ & $0.52\pm0.08$ & $0.40\pm0.08$ & 52\\
5.5 & $0.00\pm0.00$ & $0.23\pm0.05$ & $0.57\pm0.06$ & $0.21\pm0.05$ & 265\\
4.5 & $0.01\pm0.01$ & $0.48\pm0.09$ & $0.46\pm0.10$ & $0.05\pm0.03$ & 56\\
3.25 & $0.01\pm0.03$ & $0.47\pm0.13$ & $0.47\pm0.13$ & $0.05\pm0.06$ & 19\\
\enddata
\tablecomments{The $z=3.25$ row is an extrapolation below the
$4.5<z<6.5$ calibration window and shows the largest realization scatter.}
\end{deluxetable}

Figure~\ref{fig:march} and Table~\ref{tab:march} give the march. The
selected red fraction (xLRD+plusLRD) rises from $0.04\pm0.05$ at
$z\simeq9.5$ to $0.49\pm0.09$ at $z\simeq4.5$---a many-sigma monotonic
evolution---while the bLRD fraction falls from $0.61\pm0.15$ to
$0.05\pm0.03$ and the median dense photon fraction climbs from 0.40 to
$\simeq0.85$. At $z>8.5$, blue-continuum broad-line sources are the
largest predicted class and xLRDs are essentially absent.

Two checks establish that this is population physics rather than a
selection artifact. First, removing the selection weighting entirely
\emph{strengthens} the march (intrinsic red fraction
$0.004\rightarrow0.7$ over the same interval, computed for the full
catalogued population without a UV cut), and every subtype trend
retains its sign: a transparent selection gate can rescale absolute
fractions but cannot manufacture or erase a monotonic factor-of-$\sim$100
evolution that survives with no gate at all. Second, the rise from
$z\simeq9.5$ to $z\simeq4.5$ is monotonic in all four random-seed
realizations; the $z\simeq3.25$ bin, with $N_{\rm eff}\simeq19$ and the
largest seed scatter, is reported as plateau-consistent and flagged as an
extrapolation throughout.

The physical driver is simple: at high redshift most nuclear-active
sources have young reservoirs and vigorous host star formation, so
$x\lesssim1$, envelopes are unbuilt, and continua are blue; as reservoirs
feed and hosts age, $x$ grows, covering saturates, and the population
reddens. The march is the population-level shadow of individual sources
evolving blueward-to-redward through the visibility sequence.

\section{Three tests}\label{sec:tests}

\emph{(i) The slope, with existing data.} Between $z\simeq7.5$ and 4.5
the red fraction should rise at $0.10$--$0.14$ per unit redshift (range
spans seed realizations and bin pairings). A redshift split of the
existing 249-object stack sample \citep{PerezGonzalez2026} and of
spectroscopic compilations at $z\simeq5.5$ tests this without new
observations. One systematic must be controlled: fixed observed-frame
filters sample different rest-frame wavelengths across this interval, and
the surveys contributing at $z\simeq4$ are not those contributing at
$z\simeq8$. The prediction concerns the \emph{rest-frame}
$L_{5100}/L_{2500}$ classification; applying that classification
uniformly (as the stacked-sample construction already does) mitigates the
observed-frame drift, but the redshift dependence of the parent
selection's color completeness does not cancel and should be
forward-modeled or bounded before a measured slope is compared with the
prediction.

\emph{(ii) Composition at $z>8.5$.} The model predicts bLRD plurality
($0.61\pm0.15$) and xLRD absence at $z>8.5$. The falsification threshold
is quantitative: with a predicted red fraction of $0.04\pm0.05$, a sample
of ten classified $z>8.5$ LRDs containing a red majority ($\ge5$ red)
would exclude the prediction at $\ge99.9\%$ confidence even adopting the
$+1\sigma$ edge ($p_{\rm red}=0.09$; binomial
$P(\ge5/10)=1.0\times10^{-3}$), and at $2\times10^{-5}$ at the central
value.

\emph{(iii) Luminosity control.} Because the driver is the engine--host
contrast rather than luminosity, the dense-fraction--redshift
anticorrelation should persist at fixed bolometric luminosity,
distinguishing the march from a pure luminosity-selection trend.

\section{The $z\simeq9.5$ density is a completeness thermometer}
\label{sec:highz}
The same physics resolves the character of the reported high-redshift
abundance. In the forward catalogue at $z\simeq9.5$, the reported
broad-selection density is $3.18\times10^{-5}$ cMpc$^{-3}$
\citep{Rinaldi2026}; the model's selected density is
$5.0\times10^{-6}$ (the factor 6.3 deficit), while its intrinsic
$M_{\rm UV}<-18.5$ nuclear-active population totals
$1.1\times10^{-3}$---a factor $\sim$34 above the reported value (a
model-conditional reweighting; its headroom resides in sources assigned
selection probabilities of $10^{-6}$--$10^{-3}$ by the model's analytic
color gate). These sources are UV-bright ($M_{\rm UV}\simeq-19.4$),
compact, and embedded, and are removed only by the redness selection.
Rescaling that single gate within transparent limits moves the predicted
$z=9.5$ density from the factor $\sim$6 deficit to a factor $\sim$3
excess while moving the $z\le7.5$ bins by well under their errors: the
bin measures color completeness---which \citet{Rinaldi2026} describe as
essentially unconstrained at $z>6$--7---rather than formation physics.
The recent report of no evolution in the LRD number density across this
range \citep{Loiacono2026} is independently consistent with
selection-dominated demography.

Table~\ref{tab:gates} resolves an apparent paradox in these two claims:
how can the \emph{selected} $z\simeq9.5$ sample be bLRD-dominated if the
redness gate removes blue sources? Continuum subtype and the photometric
V-shape gate are distinct axes: only 0.3\% of intrinsic bLRDs pass the
gate, but the intrinsic blue population is three orders of magnitude
larger than the red one, so its thin passing tail still dominates the
selected sample while the removed 34$\times$ headroom is also blue.

\begin{deluxetable}{lccc}
\tablecaption{Per-subtype selection at $z\simeq9.5$ in the forward
catalogue ($M_{\rm UV}<-18.5$; single fiducial realization, whereas
Table~\ref{tab:march} is seed-averaged).\label{tab:gates}}
\tablehead{\colhead{subtype} & \colhead{intrinsic $n$ [cMpc$^{-3}$]} &
\colhead{selected $n$} & \colhead{pass fraction}}
\startdata
xLRD & \nodata & \nodata & \nodata\\
plusLRD & $2.4\times10^{-7}$ & $1.4\times10^{-7}$ & 0.57\\
minusLRD & $1.9\times10^{-6}$ & $1.5\times10^{-6}$ & 0.83\\
bLRD & $1.1\times10^{-3}$ & $2.8\times10^{-6}$ & 0.0026\\
\enddata
\tablecomments{No $M_{\rm UV}<-18.5$ xLRDs occur at this redshift in the
fiducial realization. The selected shares implied here (3/34/63\% for
plus/minus/bLRD) agree with the seed-averaged Table~\ref{tab:march}
within its uncertainties.}
\end{deluxetable}

Two consequences follow: no additional formation channel is required by
current number densities, and the informative measurement at $z>8.5$ is
compositional (test ii), not demographic. Injection/recovery completeness
for \emph{blue, compact} templates is the decisive survey product.

\section{Discussion}\label{sec:disc}
The march prediction is deliberately exposed. If LRD subtypes are
distinct formation routes, the mixture can be redshift-independent and
the prediction fails; the same null outcome would follow if the
dense-fraction sequence reflects orientation, so a flat measured march
would leave orientation and distinct-routes degenerate rather than
selecting between them---the discriminating outcome is an
\emph{inverted} march (red fraction rising with redshift), which would
falsify the visibility-sequence framework outright, or a confirmed march,
which would establish LRD spectral diversity as an evolutionary sequence
and calibrate the envelope-growth clock against cosmic time. The
$z\lesssim7$ portion of the prediction is insulated from the
high-redshift completeness question; the $z>8.5$ portion should be read
jointly with Section~\ref{sec:highz}.

\begin{acknowledgments}
This work is unfunded and was carried out independently. It made use of
public data products from the JWST archive and of published catalogues by
the teams cited herein. This research made use of NASA's Astrophysics Data
System and the arXiv preprint repository.
\end{acknowledgments}

\facilities{JWST (NIRSpec, NIRCam, MIRI)}

\software{Cloudy \citep{Chatzikos2023}, numpy, scipy, matplotlib, astropy,
colossus}

\section*{Data and code availability}
The extraction scripts and machine-readable march fractions with
uncertainties accompany the companion paper as an MIT-licensed code
repository at \url{https://github.com/spsingularity/lrd-covering-law}
(archival Zenodo DOI to be assigned at submission); the full
forward-catalogue grid is a separate CC-BY-4.0 data record (Zenodo DOI to
be assigned). Only public data products are used
\citep{Rinaldi2026,Ma2026,PerezGonzalez2026,Hviding2025,Sok2026}.

\begin{thebibliography}{}
\bibitem[Barro et al.(2025)]{Barro2025} Barro, G., et al. 2025,
arXiv:2512.15853
\bibitem[Chatzikos et al.(2023)]{Chatzikos2023} Chatzikos, M., et al.
2023, RMxAA, 59, 271
\bibitem[de Graaff et al.(2025)]{deGraaff2025} de Graaff, A., et al.
2025, arXiv:2503.16600
\bibitem[Greene et al.(2024)]{Greene2024} Greene, J. E., et al. 2024,
ApJ, 964, 39
\bibitem[Fu et al.(2025)]{Fu2025} Fu, S., et al. 2025, arXiv:2512.02096
\bibitem[Hviding et al.(2025)]{Hviding2025} Hviding, R. E., et al. 2025,
RUBIES broad-line data release, Zenodo, doi:10.5281/zenodo.16758549
\bibitem[Kocevski et al.(2024)]{Kocevski2024} Kocevski, D. D., et al.
2024, arXiv:2404.03576
\bibitem[Kocevski et al.(2025)]{Kocevski2025} Kocevski, D. D., et al.
2025, ApJ, 986, 126
\bibitem[Loiacono et al.(2026)]{Loiacono2026} Loiacono, F., et al. 2026,
arXiv:2606.30253
\bibitem[Ma et al.(2026)]{Ma2026} Ma, Y., et al. 2026, ApJ, 1000, 59
\bibitem[Matthee et al.(2024)]{Matthee2024} Matthee, J., et al. 2024,
ApJ, 963, 129
\bibitem[Naidu et al.(2025)]{Naidu2025} Naidu, R. P., et al. 2025,
arXiv:2503.16596
\bibitem[Pacucci \& Loeb(2025)]{PacucciLoeb2025} Pacucci, F., \& Loeb,
A. 2025, ApJL, arXiv:2506.03244
\bibitem[P\'erez-Gonz\'alez et al.(2026)]{PerezGonzalez2026}
P\'erez-Gonz\'alez, P. G., et al. 2026, arXiv:2602.20247
\bibitem[Rinaldi et al.(2026)]{Rinaldi2026} Rinaldi, P., et al. 2026,
arXiv:2604.07138
\bibitem[Sok et al.(2026)]{Sok2026} Sok, V., et al. 2026, arXiv:2606.23778
\bibitem[Tinker et al.(2008)]{Tinker2008} Tinker, J., et al. 2008, ApJ,
688, 709
\end{thebibliography}

\end{document}
